\documentclass[12pt]{article}
\usepackage{amsmath,amssymb,geometry,enumitem,booktabs}
\geometry{margin=1in}
\setlength{\parskip}{0.6em}
\setlength{\parindent}{0pt}

\begin{document}

\begin{center}
{\LARGE \textbf{ECON 300 -- Labour Economics}}\\[0.8em]
{\Large \textbf{Solutions: Final Exam Review}}\\[0.8em]
{\Large Winter 2026}
\end{center}

\vspace{0.5em}

% ============================================================
\section*{Part I: Labour Supply and Labour Force Participation}
% ============================================================

\begin{enumerate}

\item \textbf{Budget constraint with leisure and work}

Given:
\[
T=112,\qquad h=T-L,\qquad w=18,\qquad V=240
\]
where \(V\) is non-labour income.

\begin{enumerate}[label=(\alph*)]
    \item Weekly budget constraint:

    Weekly income \(C\) equals labour income plus non-labour income:
    \[
    C=w h + V
    \]
    Since \(h=T-L\),
    \[
    C=w(T-L)+V
    \]
    Substitute the numbers:
    \[
    C=18(112-L)+240
    \]
    \[
    C=2016-18L+240
    \]
    \[
    \boxed{C=2256-18L}
    \]

    \item Intercept on the income axis:

    The income intercept occurs when leisure is zero:
    \[
    L=0
    \]
    Then
    \[
    C=2256-18(0)=2256
    \]
    \[
    \boxed{2256}
    \]

    \item Intercept on the leisure axis:

    The leisure intercept occurs when income equals non-labour income plus zero work, or more directly by setting hours worked to zero:
    \[
    h=0 \Rightarrow L=T=112
    \]
    So the leisure intercept is:
    \[
    \boxed{112}
    \]

    \item Slope of the budget line:

    From
    \[
    C=2256-18L
    \]
    the slope is:
    \[
    \boxed{-18}
    \]

    Interpretation: one extra hour of leisure reduces income by \(\$18\), which is the wage.

    \item If the wage rises to \(\$24\):

    New budget line:
    \[
    C=24(112-L)+240
    \]
    \[
    C=2688-24L+240
    \]
    \[
    \boxed{C=2928-24L}
    \]

    The new income intercept is:
    \[
    C=2928 \quad \text{when } L=0
    \]
    The leisure intercept remains:
    \[
    L=112
    \]
    The slope becomes steeper:
    \[
    \boxed{-24}
    \]

    So the budget line rotates outward around the leisure intercept.
\end{enumerate}

% ------------------------------------------------------------

\item \textbf{Hours worked before and after a wage increase}

Initial leisure:
\[
L_0=80
\]
After the wage increase:
\[
L_1=76
\]
Total available time:
\[
T=112
\]

\begin{enumerate}[label=(\alph*)]
    \item Hours worked before:
    \[
    h_0=T-L_0=112-80=32
    \]

    Hours worked after:
    \[
    h_1=T-L_1=112-76=36
    \]

    \[
    \boxed{h_0=32,\qquad h_1=36}
    \]

    \item Labour supply increased because hours worked rose:
    \[
    36-32=4
    \]

    \[
    \boxed{\text{Labour supply increased by 4 hours}}
    \]

    \item Explanation:

    A wage increase has two effects:

    \textbf{Substitution effect:}
    Leisure becomes more expensive because each hour not worked now means giving up more income. This encourages the worker to consume less leisure and work more.

    \textbf{Income effect:}
    A higher wage makes the worker richer for any given number of hours worked. If leisure is a normal good, the worker may demand more leisure and work less.

    In this question, the worker reduced leisure from 80 to 76 hours, so the substitution effect dominated the income effect.
\end{enumerate}

% ------------------------------------------------------------

\item \textbf{Reservation wage and labour force participation}

A full workday is 8 hours. The worker values one hour of leisure at \(\$14\). Fixed cost of working is \(\$60\) per day.

\begin{enumerate}[label=(\alph*)]
    \item Minimum daily earnings needed to compensate for leisure forgone:

    If the worker works 8 hours, leisure forgone is 8 hours. Value of forgone leisure:
    \[
    8\times 14=112
    \]

    \[
    \boxed{\$112}
    \]

    \item Including fixed costs:

    Total daily compensation required:
    \[
    112+60=172
    \]

    \[
    \boxed{\$172}
    \]

    \item Reservation wage per hour:

    Divide total required daily earnings by 8 hours:
    \[
    \frac{172}{8}=21.5
    \]

    \[
    \boxed{\$21.50 \text{ per hour}}
    \]

    \item If market wage is \(\$21\):

    Daily earnings would be:
    \[
    21\times 8=168
    \]

    Since
    \[
    168<172
    \]
    the worker will not participate.

    \[
    \boxed{\text{No, the person will not participate}}
    \]
\end{enumerate}

% ============================================================
\section*{Part II: Income Maintenance, Welfare, and Work Incentives}
% ============================================================

\item \textbf{Income maintenance program}

Given:
\[
G=700,\qquad t=0.40,\qquad w=25
\]

\begin{enumerate}[label=(\alph*)]
    \item Disposable income while benefits are positive:

    The standard formula is:
    \[
    C=G+(1-t)wh
    \]
    Substitute:
    \[
    C=700+(1-0.40)(25)h
    \]
    \[
    C=700+0.60(25)h
    \]
    \[
    \boxed{C=700+15h}
    \]

    \item Effective net wage:

    The effective net wage is:
    \[
    (1-t)w=0.60(25)=15
    \]

    \[
    \boxed{\$15 \text{ per hour}}
    \]

    \item Break-even level of earnings:

    Benefits fall to zero when:
    \[
    G=twh
    \]
    or equivalently when gross earnings satisfy:
    \[
    wh=\frac{G}{t}
    \]
    So:
    \[
    wh=\frac{700}{0.40}=1750
    \]

    \[
    \boxed{\$1750}
    \]

    \item Corresponding hours of work:

    Since \(w=25\),
    \[
    25h=1750
    \]
    \[
    h=70
    \]

    \[
    \boxed{70 \text{ hours}}
    \]

    \item Effect of higher benefit reduction rate:

    A higher tax-back rate lowers the effective net wage:
    \[
    (1-t)w
    \]
    so workers keep less of each extra dollar earned. This weakens the incentive to work additional hours and may reduce labour supply.
\end{enumerate}

% ------------------------------------------------------------

\item \textbf{Demogrant vs. negative income tax style program}

Wage:
\[
w=20
\]

\begin{enumerate}[label=(\alph*)]
    \item Program A: demogrant of \(\$500\)

    A demogrant is a lump-sum transfer with no tax-back, so:
    \[
    C=500+20h
    \]

    \[
    \boxed{C=500+20h}
    \]

    \item Program B: guaranteed income of \(\$500\) with tax-back rate \(50\%\)

    Disposable income while benefits are positive:
    \[
    C=500+(1-0.5)(20)h
    \]
    \[
    C=500+10h
    \]

    \[
    \boxed{C=500+10h}
    \]

    \item Effective net wage under Program B:
    \[
    (1-0.5)20=10
    \]

    \[
    \boxed{\$10}
    \]

    \item Which program creates a substitution effect?

    Program A does \emph{not} change the wage rate, so it creates only an income effect.

    Program B lowers the effective wage from \(\$20\) to \(\$10\), so it creates a substitution effect.

    \[
    \boxed{\text{Program B creates a substitution effect}}
    \]

    \item Which program is more likely to reduce hours worked?

    Program A raises income without changing the wage, so it may reduce hours through an income effect.

    Program B raises income and also lowers the effective wage, so it reduces work incentives more strongly because it creates both an income effect and a substitution effect.

    \[
    \boxed{\text{Program B is more likely to reduce hours worked}}
    \]
\end{enumerate}

% ============================================================
\section*{Part III: Labour Demand in Competitive Markets}
% ============================================================

\item \textbf{Competitive firm labour demand}

Production:
\[
Q=40N-N^2
\]
Output price:
\[
P=3
\]
Wage:
\[
w=48
\]

\begin{enumerate}[label=(\alph*)]
    \item Marginal product of labour:

    Differentiate \(Q\) with respect to \(N\):
    \[
    MP_N=\frac{dQ}{dN}=40-2N
    \]

    \[
    \boxed{MP_N=40-2N}
    \]

    \item Value of marginal product:

    \[
    VMP_N=P\cdot MP_N=3(40-2N)=120-6N
    \]

    \[
    \boxed{VMP_N=120-6N}
    \]

    \item Profit-maximizing employment:

    Competitive firms hire labour until:
    \[
    VMP_N=w
    \]
    So:
    \[
    120-6N=48
    \]
    \[
    72=6N
    \]
    \[
    N=12
    \]

    \[
    \boxed{N=12}
    \]

    \item Total output at \(N=12\):

    \[
    Q=40(12)-12^2=480-144=336
    \]

    \[
    \boxed{Q=336}
    \]
\end{enumerate}

% ------------------------------------------------------------

\item \textbf{Elasticity of labour demand}

Given:
\[
\eta=-1.4,\qquad N_0=250,\qquad \%\Delta W=10\%
\]

\begin{enumerate}[label=(\alph*)]
    \item Percentage change in employment:

    By definition:
    \[
    \%\Delta N = \eta \cdot \%\Delta W
    \]
    \[
    \%\Delta N = -1.4(10\%)=-14\%
    \]

    \[
    \boxed{-14\%}
    \]

    \item New employment:

    A 14\% decrease from 250:
    \[
    N_1=250(1-0.14)=250(0.86)=215
    \]

    \[
    \boxed{215}
    \]

    \item Meaning of elastic labour demand:

    Labour demand is elastic when employment responds strongly to wage changes. Here, a 10\% wage increase causes a 14\% fall in employment, so firms are quite responsive to wage changes.
\end{enumerate}

% ------------------------------------------------------------

\item \textbf{Cost-minimizing input choice}

Given:
\[
MPP_N=20,\qquad MPP_K=10,\qquad w=40,\qquad r=15
\]

\begin{enumerate}[label=(\alph*)]
    \item 
    \[
    \frac{MPP_N}{MPP_K}=\frac{20}{10}=2
    \]

    \[
    \boxed{2}
    \]

    \item
    \[
    \frac{w}{r}=\frac{40}{15}=2.67
    \]

    \[
    \boxed{2.67}
    \]

    \item Cost minimization requires:
    \[
    \frac{MPP_N}{MPP_K}=\frac{w}{r}
    \]
    Since
    \[
    2 \neq 2.67
    \]
    the condition is not satisfied.

    \[
    \boxed{\text{No}}
    \]

    \item Since
    \[
    \frac{MPP_N}{MPP_K}<\frac{w}{r}
    \]
    labour is relatively expensive compared with its productivity contribution. The firm should use less labour and more capital.

    \[
    \boxed{\text{Use relatively less labour and more capital}}
    \]
\end{enumerate}

% ============================================================
\section*{Part IV: Monopsony and Wage Determination}
% ============================================================

\item \textbf{Monopsony equilibrium}

Supply:
\[
W=12+0.5N
\]
Demand:
\[
MRP_N=60-0.5N
\]

\begin{enumerate}[label=(\alph*)]
    \item Marginal cost of labour:

    Total labour cost:
    \[
    TLC=W\cdot N=(12+0.5N)N=12N+0.5N^2
    \]
    Differentiate:
    \[
    MCL=\frac{dTLC}{dN}=12+N
    \]

    \[
    \boxed{MCL=12+N}
    \]

    \item Monopsony employment:

    Set
    \[
    MCL=MRP_N
    \]
    \[
    12+N=60-0.5N
    \]
    \[
    1.5N=48
    \]
    \[
    N=32
    \]

    \[
    \boxed{N_M=32}
    \]

    \item Monopsony wage:

    Substitute into supply:
    \[
    W=12+0.5(32)=12+16=28
    \]

    \[
    \boxed{W_M=28}
    \]

    \item Competitive outcome:

    Under competition:
    \[
    W=MRP_N
    \]
    so
    \[
    12+0.5N=60-0.5N
    \]
    \[
    N=48
    \]
    Wage:
    \[
    W=12+0.5(48)=36
    \]

    \[
    \boxed{N_C=48,\quad W_C=36}
    \]

    \item Comparison:

    Monopsony leads to:
    \[
    N_M=32<48=N_C
    \]
    \[
    W_M=28<36=W_C
    \]

    So employment and wages are both lower under monopsony.
\end{enumerate}

% ------------------------------------------------------------

\item \textbf{Minimum wage under monopsony}

Minimum wage:
\[
W_{min}=32
\]

\begin{enumerate}[label=(\alph*)]
    \item Maximum employment available from supply curve:

    From
    \[
    32=12+0.5N
    \]
    \[
    20=0.5N
    \]
    \[
    N=40
    \]

    \[
    \boxed{40}
    \]

    \item Firm's desired employment:

    With a binding minimum wage, marginal cost is \(32\) up to \(N=40\). Set:
    \[
    32=60-0.5N
    \]
    \[
    0.5N=28
    \]
    \[
    N=56
    \]

    But only 40 workers are willing to work at that wage, so the feasible employment is:
    \[
    \boxed{40}
    \]

    \item Compare with original monopsony outcome:

    Original monopsony employment:
    \[
    32
    \]
    New employment:
    \[
    40
    \]

    \[
    \boxed{\text{Employment rises}}
    \]

    \item Explanation:

    Under monopsony, the firm hires too few workers because hiring more labour drives up wages. A minimum wage can flatten marginal labour cost over some range, allowing the firm to hire more workers and move closer to the competitive outcome.
\end{enumerate}

% ============================================================
\section*{Part V: Compensating Wage Differentials}
% ============================================================

\item \textbf{Risky vs safe jobs}

Safe job:
\[
\$28
\]
Risky job:
\[
\$34
\]

\begin{enumerate}[label=(\alph*)]
    \item Compensating differential:
    \[
    34-28=6
    \]

    \[
    \boxed{\$6 \text{ per hour}}
    \]

    \item Worker requires at least \(\$5\):

    Since
    \[
    6>5
    \]
    the risky job compensates enough.

    \[
    \boxed{\text{Choose risky job}}
    \]

    \item Worker requires at least \(\$7\):

    Since
    \[
    6<7
    \]
    the premium is not enough.

    \[
    \boxed{\text{Choose safe job}}
    \]

    \item Sorting explanation:

    Workers differ in how much compensation they require for unpleasant or risky job characteristics. More risk-tolerant workers accept risky jobs at lower premia, while more risk-averse workers choose safer jobs unless wage compensation is very high.
\end{enumerate}

% ------------------------------------------------------------

\item \textbf{Value of a statistical life}

Annual premium:
\[
\$1{,}200
\]
Increase in fatality risk:
\[
\frac{1}{20{,}000}
\]

\begin{enumerate}[label=(\alph*)]
    \item VSL:
    \[
    VSL=\frac{1200}{1/20{,}000}=1200(20{,}000)=24{,}000{,}000
    \]

    \[
    \boxed{\$24{,}000{,}000}
    \]

    \item Interpretation:

    This does not mean any specific life is ``worth'' \$24 million. It means that, aggregated across many workers, the labour market reveals a willingness to pay consistent with \$24 million for reducing one expected statistical death.
\end{enumerate}

% ============================================================
\section*{Part VI: Human Capital, Education, and Training}
% ============================================================

\item \textbf{Johny's diploma decision}

Annual direct costs:
\[
2500
\]
Forgone earnings:
\[
18{,}000
\]
Earnings without diploma:
\[
30{,}000
\]
Earnings with diploma:
\[
38{,}000
\]
Discount rate:
\[
r=0.05
\]
Working life after graduation:
\[
T=20
\]

\begin{enumerate}[label=(\alph*)]
    \item Total annual cost:
    \[
    2500+18{,}000=20{,}500
    \]

    \[
    \boxed{\$20{,}500}
    \]

    \item Present value of 3 years of costs:

    Costs are paid at the end of each year:
    \[
    PV_C=\frac{20{,}500}{1.05}+\frac{20{,}500}{(1.05)^2}+\frac{20{,}500}{(1.05)^3}
    \]

    Compute each term:
    \[
    \frac{20{,}500}{1.05}=19{,}523.81
    \]
    \[
    \frac{20{,}500}{1.1025}=18{,}594.10
    \]
    \[
    \frac{20{,}500}{1.157625}=17{,}708.67
    \]

    Sum:
    \[
    PV_C=19{,}523.81+18{,}594.10+17{,}708.67=55{,}826.58
    \]

    \[
    \boxed{PV_C\approx \$55{,}826.58}
    \]

    \item Annual earnings gain:
    \[
    38{,}000-30{,}000=8{,}000
    \]

    \[
    \boxed{\$8{,}000}
    \]

    \item Present value of benefits at time of graduation:

    Use annuity formula:
    \[
    PV_B^{(3)}=8000\left(\frac{1-(1.05)^{-20}}{0.05}\right)
    \]

    Since
    \[
    (1.05)^{-20}\approx 0.37689
    \]
    then
    \[
    \frac{1-0.37689}{0.05}=12.4622
    \]
    so
    \[
    PV_B^{(3)}=8000(12.4622)=99{,}697.60
    \]

    \[
    \boxed{PV_B^{(3)}\approx \$99{,}697.60}
    \]

    \item Discount benefits back to the present:

    Graduation is after 3 years, so:
    \[
    PV_B^{(0)}=\frac{99{,}697.60}{(1.05)^3}
    \]
    \[
    PV_B^{(0)}=\frac{99{,}697.60}{1.157625}\approx 86{,}122.02
    \]

    \[
    \boxed{PV_B^{(0)}\approx \$86{,}122.02}
    \]

    \item NPV rule:

    \[
    NPV=PV_B^{(0)}-PV_C
    \]
    \[
    NPV=86{,}122.02-55{,}826.58=30{,}295.44
    \]

    Since \(NPV>0\), Johny should invest.

    \[
    \boxed{\text{Yes, Johny should invest in the diploma}}
    \]
\end{enumerate}

% ------------------------------------------------------------

\item \textbf{General vs specific training}

\begin{enumerate}[label=(\alph*)]
    \item General training is more likely financed by the worker because it increases earnings in many firms, so the worker captures most of the benefit.

    \item Specific training is more likely financed by the firm because the productivity gain is mainly valuable inside the current firm.

    \item Why it matters for wages:

    During general training, trainees may accept lower wages because they will gain future transferable benefits. Under specific training, firms often share more of the cost because the future gains are not fully portable.
\end{enumerate}

% ============================================================
\section*{Part VII: Discrimination and Oaxaca Decomposition}
% ============================================================

\item \textbf{Oaxaca decomposition}

Men:
\[
\ln W_m=1.10+0.08S+0.05EXP
\]
Women:
\[
\ln W_f=0.90+0.07S+0.04EXP
\]

Means:
\[
\bar S_m=16,\quad \bar{EXP}_m=12
\]
\[
\bar S_f=15,\quad \bar{EXP}_f=10
\]

\begin{enumerate}[label=(\alph*)]
    \item Predicted log earnings for men:
    \[
    \ln W_m=1.10+0.08(16)+0.05(12)
    \]
    \[
    =1.10+1.28+0.60=2.98
    \]

    \[
    \boxed{2.98}
    \]

    \item Predicted log earnings for women:
    \[
    \ln W_f=0.90+0.07(15)+0.04(10)
    \]
    \[
    =0.90+1.05+0.40=2.35
    \]

    \[
    \boxed{2.35}
    \]

    \item Total log wage gap:
    \[
    2.98-2.35=0.63
    \]

    \[
    \boxed{0.63}
    \]

    \item Explained portion using male structure:

    Difference in schooling:
    \[
    16-15=1
    \]
    Difference in experience:
    \[
    12-10=2
    \]

    Multiply by male coefficients:
    \[
    (1)(0.08)+(2)(0.05)=0.08+0.10=0.18
    \]

    \[
    \boxed{0.18}
    \]

    \item Unexplained portion:
    \[
    0.63-0.18=0.45
    \]

    \[
    \boxed{0.45}
    \]

    \item Interpretation:

    Of the total log wage gap of 0.63, 0.18 is explained by measured differences in schooling and experience. The remaining 0.45 is unexplained and may reflect discrimination or omitted variables.
\end{enumerate}

% ============================================================
\section*{Part VIII: Immigration}
% ============================================================

\item \textbf{Immigration categories}

Total:
\[
360{,}000
\]

\begin{enumerate}[label=(\alph*)]
    \item Shares:

    Economic:
    \[
    \frac{180{,}000}{360{,}000}\times 100=50\%
    \]

    Family:
    \[
    \frac{108{,}000}{360{,}000}\times 100=30\%
    \]

    Refugees:
    \[
    \frac{72{,}000}{360{,}000}\times 100=20\%
    \]

    \[
    \boxed{50\%,\ 30\%,\ 20\%}
    \]

    \item Economic immigrants rise by 15\%:
    \[
    180{,}000(1.15)=207{,}000
    \]

    New total:
    \[
    207{,}000+108{,}000+72{,}000=387{,}000
    \]

    \[
    \boxed{387{,}000}
    \]

    \item New economic share:
    \[
    \frac{207{,}000}{387{,}000}\times 100\approx 53.49\%
    \]

    \[
    \boxed{53.49\%}
    \]
\end{enumerate}

% ------------------------------------------------------------

\item \textbf{Migration decision}

Earnings:
\[
54{,}000-48{,}000=6{,}000
\]
So the annual advantage is:
\[
\boxed{\$6{,}000}
\]

\begin{enumerate}[label=(\alph*)]
    \item Annual advantage:
    \[
    \boxed{\$6{,}000}
    \]

    \item Present value:

    \[
    PV=6000\left(\frac{1-(1.04)^{-25}}{0.04}\right)
    \]

    Since
    \[
    (1.04)^{-25}\approx 0.3751
    \]
    then
    \[
    \frac{1-0.3751}{0.04}=15.6225
    \]
    so
    \[
    PV=6000(15.6225)=93{,}735
    \]

    \[
    \boxed{\$93{,}735}
    \]

    \item Subtract moving cost:
    \[
    93{,}735-20{,}000=73{,}735
    \]

    Since this is positive, the migrant should move.

    \[
    \boxed{\text{Yes, the migrant should move}}
    \]
\end{enumerate}

% ============================================================
\section*{Part IX: Unemployment Measurement and Dynamics}
% ============================================================

\item \textbf{Basic labour market indicators}

Population:
\[
8{,}000{,}000
\]
Employed:
\[
5{,}360{,}000
\]
Unemployed:
\[
340{,}000
\]

\begin{enumerate}[label=(\alph*)]
    \item Labour force:
    \[
    5{,}360{,}000+340{,}000=5{,}700{,}000
    \]

    \[
    \boxed{5{,}700{,}000}
    \]

    \item Unemployment rate:
    \[
    \frac{340{,}000}{5{,}700{,}000}\times 100=5.96\%
    \]

    \[
    \boxed{5.96\%}
    \]

    \item Labour force participation rate:
    \[
    \frac{5{,}700{,}000}{8{,}000{,}000}\times 100=71.25\%
    \]

    \[
    \boxed{71.25\%}
    \]

    \item Employment rate:
    \[
    \frac{5{,}360{,}000}{8{,}000{,}000}\times 100=67.00\%
    \]

    \[
    \boxed{67.00\%}
    \]
\end{enumerate}

% ------------------------------------------------------------

\item \textbf{Changes in unemployment}

Labour force:
\[
2{,}500{,}000
\]
Unemployment rate:
\[
8.4\%
\]

\begin{enumerate}[label=(\alph*)]
    \item Unemployed:
    \[
    0.084(2{,}500{,}000)=210{,}000
    \]

    \[
    \boxed{210{,}000}
    \]

    \item Employed:
    \[
    2{,}500{,}000-210{,}000=2{,}290{,}000
    \]

    \[
    \boxed{2{,}290{,}000}
    \]

    \item New unemployed workers:

    35,000 unemployed find jobs:
    \[
    210{,}000-35{,}000=175{,}000
    \]

    20,000 employed workers lose jobs and search:
    \[
    175{,}000+20{,}000=195{,}000
    \]

    \[
    \boxed{195{,}000}
    \]

    \item New unemployment rate:
    \[
    \frac{195{,}000}{2{,}500{,}000}\times 100=7.80\%
    \]

    \[
    \boxed{7.80\%}
    \]
\end{enumerate}

% ------------------------------------------------------------

\item \textbf{Incidence and duration}

\begin{enumerate}[label=(\alph*)]
    \item Region A:
    \[
    UR=3\%\times 2.5=7.5\%
    \]

    Region B:
    \[
    UR=1.5\%\times 5=7.5\%
    \]

    \[
    \boxed{UR_A=7.5\%,\quad UR_B=7.5\%}
    \]

    \item Higher incidence:
    \[
    3\%>1.5\%
    \]
    \[
    \boxed{\text{Region A}}
    \]

    \item Longer duration:
    \[
    5>2.5
    \]
    \[
    \boxed{\text{Region B}}
    \]

    \item Policy explanation:

    The same unemployment rate can come from many short unemployment spells or fewer long spells. Region A's problem is frequent entry into unemployment. Region B's problem is prolonged joblessness. These require different policies.
\end{enumerate}



\end{enumerate}

\end{document}